\documentclass[aspectratio=169,10pt,english]{beamer}

\input{../../../_preambulo_beamer.tex}
\input{../../../_comandos_beamer.tex}

\selectlanguage{english}

\title{MA1001B - Statistical Modeling for Decision Making}
\subtitle{Lesson 08: Probability Rules and Contingency Tables}
\author{}
\institute{Tecnol\'ogico de Monterrey}
\date{}

\begin{document}

\begin{frame}[plain]
  \vspace{1.2cm}
  {\Large\bfseries MA1001B - Statistical Modeling for Decision Making\par}
  \vspace{0.7cm}
  {\large Lesson 08: Probability Rules and Contingency Tables\par}
  \vspace{0.7cm}
  {\color{TecRojo}\rule{\textwidth}{1.2pt}}
  \vspace{0.8cm}

  Extensive material for independent study\\[0.3cm]
  Synthetic case of 400 orders\\[0.3cm]
  Tecnol\'ogico de Monterrey
\end{frame}

\begin{frame}{The Probability Question}
  \begin{alertblock}{Guiding question}
    How can one contingency table support correct AND, OR, and NOT
    probabilities without double-counting or assuming independence?
  \end{alertblock}

  \vspace{0.4em}
  By the end of this lesson, you can:
  \begin{itemize}
    \item read marginal and joint probabilities from a two-way table;
    \item apply complement and general addition rules;
    \item identify the overlap in an OR probability;
    \item explain why multiplying marginal probabilities requires evidence of
      independence.
  \end{itemize}

  \vfill
  \scriptsize All order records and operating processes are synthetic.
\end{frame}

\begin{frame}{A Contingency Table Organizes Joint Outcomes}
  \small
  Let \(M\) be manual review and \(L\) be late delivery.

  \begin{center}
    \begin{tabular}{lrrr}
      \toprule
      \textbf{Routing process} & \textbf{Late} & \textbf{On time} & \textbf{Total} \\
      \midrule
      Manual review & 48 & 72 & 120 \\
      Automated & 32 & 248 & 280 \\
      \midrule
      Total & 80 & 320 & 400 \\
      \bottomrule
    \end{tabular}
  \end{center}

  \vspace{0.6em}
  \begin{columns}[T]
    \begin{column}{0.48\textwidth}
      \textbf{Marginal probabilities}
      \[
        P(M)=\frac{120}{400}=0.30,\quad
        P(L)=\frac{80}{400}=0.20.
      \]
    \end{column}
    \begin{column}{0.48\textwidth}
      \textbf{Joint probability}
      \[
        P(M\cap L)=\frac{48}{400}=0.12.
      \]
    \end{column}
  \end{columns}
\end{frame}

\begin{frame}{Step 1: Read Cells and Margins Before Applying Rules}
  \begin{columns}[T]
    \begin{column}{0.36\textwidth}
      \small
      \begin{block}{Verified output}
        Synthetic orders: 400\\
        \(P(M)=0.30\)\\
        \(P(L)=0.20\)\\
        \(P(M\cap L)=0.12\)
      \end{block}

      A cell represents two conditions at once. Row and column totals represent
      one condition regardless of the other variable.
    \end{column}
    \begin{column}{0.64\textwidth}
      \centering
      \includegraphics[width=0.88\textwidth]{../figures/probability_01_table.png}
    \end{column}
  \end{columns}
\end{frame}

\begin{frame}{Complement and Addition Rules}
  \small
  The complement rule accounts for everything outside an event:
  \[
    P(M^c)=1-P(M)=1-0.30=0.70.
  \]

  The general addition rule removes the shared outcomes counted twice:
  \[
    P(M\cup L)=P(M)+P(L)-P(M\cap L).
  \]
  \[
    P(M\cup L)=0.30+0.20-0.12=0.38.
  \]

  \begin{alertblock}{Common trap}
    The naive sum 0.50 includes the 48 manual-and-late orders twice. The simple
    sum applies only when the events are mutually exclusive.
  \end{alertblock}
\end{frame}

\begin{frame}{Step 2: The Addition Rule Matches Direct Counting}
  \begin{columns}[T]
    \begin{column}{0.35\textwidth}
      \small
      \begin{block}{Verified comparison}
        Naive sum: 0.50\\
        Overlap: 0.12\\
        Corrected union: 0.38\\
        Direct table count: 0.38
      \end{block}

      Directly counting all orders in \(M\cup L\) gives
      \[
        \frac{152}{400}=0.38.
      \]
    \end{column}
    \begin{column}{0.65\textwidth}
      \centering
      \includegraphics[width=\textwidth]{../figures/probability_02_addition.png}
    \end{column}
  \end{columns}
\end{frame}

\begin{frame}{Step 3: The Multiplication Shortcut Has a Condition}
  \begin{columns}[T]
    \begin{column}{0.38\textwidth}
      \small
      General rule:
      \[
        P(M\cap L)=P(M)P(L\mid M).
      \]
      The table gives \(P(L\mid M)=48/120=0.40\), so
      \[
        0.30(0.40)=0.12.
      \]

      Multiplying marginals gives \(0.30(0.20)=0.06\), which contradicts the
      observed intersection.

      \textbf{Limit:} independence has not been justified.
    \end{column}
    \begin{column}{0.62\textwidth}
      \centering
      \includegraphics[width=\textwidth]{../figures/probability_03_limit.png}
    \end{column}
  \end{columns}
\end{frame}

\begin{frame}{From Records to a Table}
  \small
  Each of the 400 synthetic orders has one routing category and one delivery
  outcome. \texttt{pd.crosstab} counts each combination once.
  \begin{center}
    \begin{tabular}{lll}
      \toprule
      Selection & Count & Probability\\
      \midrule
      Manual review & \(48+72=120\) & \(0.30\)\\
      Late & \(48+32=80\) & \(0.20\)\\
      Manual and late & \(48\) & \(0.12\)\\
      Automated and on time & \(248\) & \(0.62\)\\
      \bottomrule
    \end{tabular}
  \end{center}
  A cell is an intersection; a margin ignores the other variable.
  The denominator is 400 because each row represents one order.
\end{frame}

\begin{frame}{The Union Removes a Duplicate Count}
  \small
  The 120 manual-review orders include 48 late orders. The 80 late orders
  include those same 48.
  \[
    |M\cup L|=120+80-48=152,\qquad
    P(M\cup L)=\frac{152}{400}=0.38.
  \]
  \begin{block}{Independent check}
    Only the 248 automated, on-time orders lie outside \(M\cup L\).
    By the complement, \(P(M\cup L)=1-248/400=0.38\).
  \end{block}
  Simple addition works only for an empty intersection. Here
  \(0.30+0.20=0.50\) exceeds the true union by \(0.12\).
\end{frame}

\begin{frame}{Conditioning Changes the Denominator}
  \small
  Among 120 manual-review orders, 48 were late:
  \[
    P(L\mid M)=\frac{48}{120}=0.40.
  \]
  Among 80 late orders, 48 were manually reviewed:
  \[
    P(M\mid L)=\frac{48}{80}=0.60.
  \]
  The intersection stays the same, but the reference population changes.
  The general rule matches the table cell:
  \[
    P(M\cap L)=P(M)P(L\mid M)=0.30(0.40)=0.12.
  \]
  Replacing the conditional with marginal \(P(L)=0.20\) gives \(0.06\).
\end{frame}

\begin{frame}{How the Scripts Build Evidence}
  \small
  \begin{enumerate}
    \item \texttt{build\_operations\_data()} expands four fixed counts into
      400 synthetic rows. Seed 42 controls their order only.
    \item \texttt{pd.crosstab} recovers the four cells and checks the total.
    \item \texttt{table\_probabilities()} divides margins and cells by 400.
    \item Step 2 compares the addition rule with a direct count.
    \item Step 3 compares the marginal product with the general rule.
  \end{enumerate}
  All three programs run independently and save figures without opening
  windows. The study notebook integrates them with editable practice.
\end{frame}

\begin{frame}{Worked Example: Automated or Late}
  \small
  Let \(A\) mean \emph{automated}. The table gives
  \(P(A)=280/400=0.70\), \(P(L)=80/400=0.20\), and
  \(P(A\cap L)=32/400=0.08\).
  \[
    P(A\cup L)=0.70+0.20-0.08=0.82.
  \]
  Direct count: \(280+80-32=328\), so \(328/400=0.82\).
  \begin{block}{Practice}
    Find \(P(\text{on time}\cup M)\), then check it through the
    complementary cell.
  \end{block}
\end{frame}

\begin{frame}{Scope of This Table}
  \small
  These frequencies belong to a fabricated population of 400 orders.
  They do not represent a company or estimate an external population rate.
  The association between manual review and lateness does not establish
  causality: there is no random assignment or control for other factors.

  \begin{block}{Conceptual bridge}
    L09 develops the reduced sample space created by conditioning
    and formal checks of independence.
  \end{block}
\end{frame}

\begin{frame}{Concept Check}
  Use the 400-order table from this lesson.
  \begin{enumerate}
    \item Find \(P(\text{On time})\).
    \item Find \(P(\text{Automated OR Late})\) using the addition rule.
    \item Why is \(P(M)+P(L)=0.50\) not the probability of \(M\cup L\)?
    \item What evidence shows that \(P(M)P(L)\) is not valid for this
      intersection?
  \end{enumerate}

  \vspace{0.4em}
  \begin{block}{Connection to Lesson 09}
    The next lesson develops conditional probability and independence: how a
    reduced sample space changes probability and how independence is tested.
  \end{block}

  \vfill
  \scriptsize Source: OpenStax, \textit{Introductory Business Statistics 2e},
  Chapter 3, Sections 3.1 and 3.3. CC BY-NC-SA 4.0.
\end{frame}

\appendix



\end{document}
